\documentclass[11pt]{article}

\usepackage[margin=1in]{geometry}
\usepackage{amsmath,amssymb}
\usepackage{booktabs}
\usepackage{siunitx}
\usepackage{hyperref}
\usepackage{graphicx}
\usepackage{xcolor}

\title{SIIMPL Training Data for the GVP--EGNN Reconstruction Model\\
\large Channeling-aware sampling, lattice symmetry, and generation parameters}
\author{Walsworth Group --- ML-Driven Damage Track Reconstruction}
\date{\today}

\begin{document}
\maketitle

\section{Motivation}

The v2 GVP--EGNN reconstruction model was trained on \num{197000} SRIM-generated
carbon recoil tracks in diamond. SRIM treats the target as an amorphous medium
with isotropic scattering, which permits the use of full $\mathrm{SO}(3)$
augmentation during training but omits the crystalline structure of diamond and
in particular \emph{channeling}: the deflection-suppressed transport of recoils
that travel within a small cone of low-index lattice axes. To extend the v2
result to a physically realistic crystalline target, the training and evaluation
data must come from SIIMPL, which propagates recoils against the actual diamond
lattice and reproduces channeling. This document specifies (i) the channeling
physics that informs sampling, (ii) the training/evaluation dataset structure,
and (iii) the SIIMPL generation parameters.

\section{Channeling physics relevant to sampling}

\subsection{Critical angle (Lindhard / Gemmell)}

The continuum-string approximation~\cite{Lindhard1965} gives the
\emph{Lindhard critical angle} for axial channeling in the high-energy
regime~\cite[Eq.~(2.12)]{Gemmell1974}:
\begin{equation}
  \psi_c \;=\; \sqrt{\frac{2 Z_1 Z_2 e^2}{E\,d}},
  \label{eq:psi_c}
\end{equation}
where $Z_1, Z_2$ are the projectile and target atomic numbers, $E$ is the
projectile kinetic energy, $d$ is the atomic spacing along the channel row, and
$e^2 = 14.4~\mathrm{eV\cdot\text{\AA}}$ is the squared electronic
charge~\cite[Glossary]{Gemmell1974}. For carbon-on-carbon recoils
($Z_1 = Z_2 = 6$),
\begin{equation}
  2 Z_1 Z_2 e^2 \;=\; 2 \cdot 6 \cdot 6 \cdot 14.4 \;=\; 1036.8~\mathrm{eV\cdot\text{\AA}}.
\end{equation}

The validity boundary between the high-energy form~\eqref{eq:psi_c} and the
low-energy $\psi_2$ form~\cite[Eq.~(2.28)]{Gemmell1974} is set
by~\cite[Eq.~(2.29)]{Gemmell1974}
\begin{equation}
  E' \;=\; \frac{2 Z_1 Z_2 e^2}{a^2}.
\end{equation}
For diamond ($a = \SI{3.567}{\angstrom}$), $E' \approx \SI{82}{eV}$, comfortably
below our minimum recoil energy of \SI{1}{keV}. The high-energy form is
therefore valid across the entire $[1,105]~\mathrm{keV}$ window and no
low-energy correction is required.

\subsection{Channeling directions in diamond}

Diamond has cubic symmetry (space group $Fd\bar{3}m$) with lattice constant
$a = \SI{3.567}{\angstrom}$~\cite{Kittel}. The three families of low-index
axial channels are determined entirely by cube symmetry and are listed in
Table~\ref{tab:axes}. The atomic spacings $d$ along each row are derived
geometrically from $a$ and the two-atom basis of the diamond structure
(atoms at $(0,0,0)$ and $(a/4, a/4, a/4)$).

\begin{table}[h]
  \centering
  \begin{tabular}{lccc}
    \toprule
    Axis family & Direction & $d$ formula & $d$ (\AA) \\
    \midrule
    $\langle 110\rangle$ & face diagonal & $a/\sqrt{2}$        & 2.52 \\
    $\langle 111\rangle$ & body diagonal & $\sqrt{3}\,a/2$ (mean of dimer spacings) & 3.09 \\
    $\langle 100\rangle$ & cube edge    & $a$                  & 3.57 \\
    \bottomrule
  \end{tabular}
  \caption{Low-index channeling directions in diamond. The $\langle 111\rangle$
    row has alternating short ($\sqrt{3}\,a/4 \approx 1.54$~\AA) and long
    ($3\sqrt{3}\,a/4 \approx 4.63$~\AA) spacings owing to the two-atom basis;
    the appropriate $d$ for the continuum-string critical-angle expression is
    the mean spacing $\sqrt{3}\,a/2$, as used in the experimental
    treatment of Derry et~al.~\cite{Derry1982} (see appendix therein for the
    planar analogue).}
  \label{tab:axes}
\end{table}

The number of distinct equivalent directions per family (unsigned): three
$\langle 100\rangle$, six $\langle 110\rangle$, four $\langle 111\rangle$.

\subsection{Critical angle values}

Evaluating Eq.~\eqref{eq:psi_c} with the $d$ values of Table~\ref{tab:axes}
gives the critical angles in Table~\ref{tab:psi}.

\begin{table}[h]
  \centering
  \begin{tabular}{lccc}
    \toprule
    & $\psi_c$ at \SI{1}{keV} & $\psi_c$ at \SI{10}{keV} & $\psi_c$ at \SI{100}{keV} \\
    \midrule
    $\langle 110\rangle$ ($d=2.52$~\AA) & $36.6^\circ$ & $11.6^\circ$ & $3.7^\circ$ \\
    $\langle 111\rangle$ ($d=3.09$~\AA) & $33.1^\circ$ & $10.5^\circ$ & $3.3^\circ$ \\
    $\langle 100\rangle$ ($d=3.57$~\AA) & $30.8^\circ$ & $9.8^\circ$  & $3.1^\circ$ \\
    \bottomrule
  \end{tabular}
  \caption{Critical channeling angles $\psi_c$ for C-on-C recoils in diamond,
    computed from Eq.~\eqref{eq:psi_c} with the row spacings of
    Table~\ref{tab:axes}.}
  \label{tab:psi}
\end{table}

The strict theoretical ordering $\langle 110\rangle > \langle 111\rangle >
\langle 100\rangle$ agrees with the experimental measurements of Derry
et~al.~\cite{Derry1982} on natural diamond. At \SI{1}{MeV} protons and room
temperature, Derry et~al.\ report measured half-widths $\psi_{1/2} = 0.55^\circ$,
$0.49^\circ$, $0.42^\circ$ for $\langle 110\rangle$, $\langle 111\rangle$,
$\langle 100\rangle$ respectively (their Table~II and Fig.~7).

The raw Lindhard $\psi_c$ values of Table~\ref{tab:psi} systematically
overpredict the experimentally measured half-widths. Derry
et~al.~\cite{Derry1982}, following Barrett, find good agreement between
experiment and the corrected angle
\begin{equation}
  \psi_{1/2} \;=\; k\,\psi_c, \qquad k \approx 0.83~\text{(axes)},
\end{equation}
where $k$ absorbs the effect of thermal vibrations and the Thomas--Fermi
screening distance. For Pool~B sampling within $2\psi_c$ cones (Sec.~3.1),
the factor of two amply absorbs the $\sim 17\%$ Barrett correction, so the
correction is presentational rather than affecting the sampler.

\section{Training dataset}

\subsection{Energy binning}

Both pools draw recoil energies from
$\mathrm{LogUniform}[\SI{1}{keV}, \SI{105}{keV}]$, matching the v2 prior.
For stratified analysis and for sizing Pool~B (Sec.~3.3) we partition the energy
range into three equal log-bins, listed in Table~\ref{tab:ebins}.

\begin{table}[h]
  \centering
  \begin{tabular}{lcccc}
    \toprule
    Bin   & Range (keV)        & Midpoint (log) & Pool A share \\
    \midrule
    Low   & $[1.00,\,4.66)$    & 2.15           & $\approx 33\%$ \\
    Mid   & $[4.66,\,21.7)$    & 10.0           & $\approx 33\%$ \\
    High  & $[21.7,\,105]$     & 46.4           & $\approx 33\%$ \\
    \bottomrule
  \end{tabular}
  \caption{Three equal log-spaced energy bins on $[1, 105]~\mathrm{keV}$. The
    bin edges are placed at $10^{0},\,10^{0.674},\,10^{1.348},\,10^{2.02}$ so
    that under $\mathrm{LogUniform}$ sampling each bin receives an equal share
    of tracks.}
  \label{tab:ebins}
\end{table}

Combined with the three axis families of Table~\ref{tab:axes}, the analysis
grid has $3 \times 3 = 9$ (axis $\times$ energy) bins.

\subsection{Pool A --- uniform on $S^2$}

Pool~A is generated first, with \num{200000} training tracks and \num{40000}
evaluation tracks. Per-track sampling:
\begin{enumerate}
  \item Draw $E \sim \mathrm{LogUniform}[\SI{1}{keV}, \SI{105}{keV}]$.
  \item Draw $g_1, g_2, g_3 \sim \mathcal{N}(0,1)$ independently and set
        $\hat{\mathbf{v}} = (g_1, g_2, g_3)/\|\mathbf{g}\|$. The resulting unit
        vector is uniformly distributed on the unit sphere.
\end{enumerate}

\subsection{Pool B --- deficit-driven channeling fill}

Pool~A's natural channeling coverage is sufficient at low and mid energies but
thin in the high-$E$ bin (where $\psi_c$ is small). Rather than picking a-priori
axis weights, we size Pool~B \emph{after} Pool~A is generated, filling each of
the nine (axis $\times$ energy) bins up to a target count.

\paragraph{Procedure.}
\begin{enumerate}
  \item Generate Pool~A.
  \item For each (axis family $\langle hkl\rangle$, energy bin $b$), count the
        Pool~A tracks whose direction lies within $\psi_c(E,d_{\langle hkl\rangle})$
        of any equivalent direction in the family. Call this $N_A^{\langle hkl\rangle, b}$.
  \item Set per-bin Pool~B target
        $N_B^{\langle hkl\rangle, b} \;=\; \max\!\left(0,\; N_{\text{target}}
        - N_A^{\langle hkl\rangle, b}\right)$ with
        $N_{\text{target}} = \num{3000}$.
  \item Generate Pool~B tracks per bin: draw $E$ uniformly within the bin,
        pick one of the family's equivalent directions uniformly, and sample
        $\hat{\mathbf{v}}$ inside a cone of half-angle $2\psi_c$ about that
        direction. The $2\psi_c$ cone exceeds the strict channeling boundary so
        the network is exposed to the transition zone in addition to deeply
        channeled trajectories.
\end{enumerate}

\paragraph{Target rationale.} $N_{\text{target}} = \num{3000}$ channeled tracks
per bin gives a $\sim 10\times$ safety margin over the statistical-precision
floor required for $\pm 0.3^\circ$ uncertainty on a per-bin median angular
error, and is comfortably above the morphology-learning threshold for the
conditional flow ($\sim \num{1000}$). The resulting Pool~B size is expected to
be $\sim \num{50000}$ tracks (mostly in the high-$E$ bin and the
under-represented $\langle 100\rangle$ family); the exact size is determined by
the measurement.

\subsection{Total dataset}

\begin{center}
  \begin{tabular}{lcc}
    \toprule
    & Train (target) & Eval (target) \\
    \midrule
    Pool A (uniform-$S^2$)             & \num{200000} & \num{40000} \\
    Pool B (deficit-filled)            & $\sim$\num{50000} & \num{10000} \\
    \midrule
    \textbf{Total}                     & $\sim$\textbf{\num{250000}} & \textbf{\num{50000}} \\
    \bottomrule
  \end{tabular}
\end{center}
The $\sim$\num{250000} training-track scale is chosen as $\approx 25\%$ above
the v2 SRIM corpus (\num{197000} tracks) to absorb the additional physical
complexity introduced by channeling, while remaining well within the local
SIIMPL generation budget.

\paragraph{Learning-curve validation.} After training, we retrain the model on
\num{50000}, \num{100000}, \num{200000}, and \num{250000} subsets and plot eval
median axis-error against training-set size on a log-$x$ axis. A flattening of
the curve by \num{250000} confirms the chosen scale is not data-limited; if the
curve is still decreasing, an additional \num{50000} tracks are generated and
the procedure is rerun. This cost is $\sim 5$ additional A100-hours and
converts the dataset-size choice from an inherited default into an empirical
claim.

\section{Crystal-orientation handling}

\subsection{Generation-time: fixed crystal frame}

SIIMPL is run with \texttt{rotate\_crystal\_mode=False}; the lab frame coincides
with the crystal frame for every event. Per-event crystal-orientation
randomization is \emph{not} used because it would average over the channeling
signal we generate the data to capture.

\subsection{Training-time: $\mathrm{O}_h$ augmentation}

The v2 training script applies random $\mathrm{SO}(3)$ rotations to each input
cloud and its label direction. This is appropriate for an amorphous SRIM
target (isotropic) but incorrect for SIIMPL: a generic $\mathrm{SO}(3)$
rotation maps a channeling direction such as $\langle 110\rangle$ to a generic
direction that the diamond lattice does not have, destroying the channeling
signal.

The correct symmetry group is the octahedral group $\mathrm{O}_h$, the 48
rotations and improper rotations that map the diamond lattice to itself:
24 proper rotations of the cube plus 24 obtained by composition with
inversion. We precompute the 48 matrices once as a tensor of shape
$(48,3,3)$. At each training step, an index is sampled uniformly from
$\{1,\dots,48\}$ and the corresponding matrix is applied to both the
point-cloud coordinates and the label direction. The model thus learns
$\mathrm{O}_h$ invariance while remaining sensitive to channeling.

\subsection{Inference-time: $\mathrm{O}_h$ averaging}

For headline-metric reporting we apply $\mathrm{O}_h$ averaging: the input
cloud is rotated through all 48 group elements, the model is evaluated on each
rotation, and posteriors are averaged. The result is exactly $\mathrm{O}_h$-
invariant with reduced variance. Single-shot inference is reported separately
for the throughput number.

\section{Two-stage training protocol}

The v2 two-stage protocol is retained with one modification (the Stage-2 pool
restriction below):
\begin{description}
  \item[Stage 1 (joint, 180 epochs).] The full \num{200000}-track training set
    (Pool A $\cup$ Pool B) is used to train the GVP--EGNN backbone, the direction
    and energy heads, and the conditional normalizing flow simultaneously.
    The backbone sees both uniform and channeled morphologies.
  \item[Stage 2 (flow-only, 60 epochs).] The backbone and heads are frozen;
    only the conditional flow is fine-tuned, and only on Pool A
    (\num{160000} uniform-on-$S^2$ tracks). The flow's implicit prior is thus
    uniform-on-$S^2$, and posterior credible regions are calibrated against
    the prior under which results are reported.
\end{description}

\section{SIIMPL generation parameters}

\begin{center}
  \begin{tabular}{ll}
    \toprule
    Parameter & Value \\
    \midrule
    Target material               & Single-crystal diamond \\
    Lattice constant $a$          & \SI{3.567}{\angstrom} \\
    Displacement energy $E_d$     & \SI{43}{eV} (locked) \\
    Projectile species            & C ($Z=6$, $A=12$) \\
    Energy distribution           & $\mathrm{LogUniform}[\SI{1}{keV}, \SI{105}{keV}]$ \\
    Direction sampling            & Pool-dependent (\S 3.1) \\
    Crystal-frame rotation        & Disabled \\
    Entry geometry                & Default random implant point \\
    Tracking mode                 & Full cascade, all vacancies \\
    Output format                 & SIIMPL native CSV $\to$ v2 schema via $\sim 20$-line parser \\
    Train/eval RNG seed split     & Disjoint master seeds \\
    \bottomrule
  \end{tabular}
\end{center}

The output schema after the parser matches v2:
\begin{center}
  \texttt{x, y, z, ion\_number, energy\_keV, theta\_deg,}\\
  \texttt{target\_vx, target\_vy, target\_vz}
\end{center}
with per-vacancy rows and per-ion labels broadcast across them.

\section{Summary}

\begin{itemize}
  \item Channeling directions in diamond ($\langle 100\rangle$,
        $\langle 110\rangle$, $\langle 111\rangle$) are forced by cubic
        symmetry~\cite{Kittel}; their associated row spacings $d$ follow from
        the lattice constant and the two-atom basis.
  \item The critical angle $\psi_c$ follows Eq.~\eqref{eq:psi_c}, valid across
        $1$--$\SI{105}{keV}$ in diamond (validity boundary
        $E' \approx \SI{82}{eV}$). The ordering
        $\psi_c^{\langle 110\rangle} > \psi_c^{\langle 111\rangle} >
        \psi_c^{\langle 100\rangle}$ agrees with the experimental
        measurements of Derry et~al.~\cite{Derry1982}.
  \item Training data: $\sim$\num{250000} tracks. Pool A (\num{200000})
        uniform-on-$S^2$ is generated first; Pool B ($\sim$\num{50000}) is
        sized post-hoc to bring each of the nine
        (axis $\times$ energy) bins up to a target of \num{3000} channeled
        tracks, with Pool B samples drawn from $2\psi_c$ cones around the
        chosen axes.
  \item Crystal-orientation invariance handled via the 48-element octahedral
        group $\mathrm{O}_h$ at training and inference, replacing v2's
        $\mathrm{SO}(3)$ augmentation.
  \item Stage~1 trains on the combined pool; Stage~2 fine-tunes the flow on
        Pool A only, keeping the posterior prior uniform-on-$S^2$.
  \item Evaluation: \num{50000} tracks (\num{40000} uniform + \num{10000}
        stratified) with full SBI calibration (SBC at $N=5000$, TARP,
        per-parameter ECE), headline metrics by energy decade, geometric
        baselines, and a zero-shot v2-on-SIIMPL cross-domain reference number.
\end{itemize}

\begin{thebibliography}{9}

\bibitem{Lindhard1965}
  J.~Lindhard,
  ``Influence of crystal lattice on motion of energetic charged particles,''
  \emph{Mat.\ Fys.\ Medd.\ Dan.\ Vid.\ Selsk.} \textbf{34}, no.~14 (1965).

\bibitem{Gemmell1974}
  D.~S.~Gemmell,
  ``Channeling and related effects in the motion of charged particles through crystals,''
  \emph{Rev.\ Mod.\ Phys.} \textbf{46}, 129 (1974).

\bibitem{Kittel}
  C.~Kittel,
  \emph{Introduction to Solid State Physics}, 8th ed.\ (Wiley, 2005), Ch.~1.

\bibitem{Derry1982}
  T.~E.~Derry, R.~W.~Fearick, and J.~P.~F.~Sellschop,
  ``Ion channeling in natural diamond. II.\ Critical angles,''
  \emph{Phys.\ Rev.\ B} \textbf{26}, 17 (1982).

\end{thebibliography}

\end{document}
